% STATUS: rewritten 2026-11-05
% STATUS: extended 2026-11-07 (Elduque 2015 read in full for the Okubo
% derivation theorems)
% Sources read in full: Baez 2002 (pp. 148, 159), Davenport 1972 (p. 565),
% Filippov 2000 (Theorems 2 and 3); Elduque 2015 (pp. 3--9).
\section{Alternative and Malcev algebras}\label{sec:nonassociative:alternative}

The previous section built the Albert algebra from the Cayley algebra % claim-ok: link sentence recalling Section 1, no external source
$\mathbb{K}$. An algebra $A$ is
\emph{alternative} if the subalgebra generated by any two elements is
associative \cite[p.~148]{Baez2002Octonions}. This is the same as asking
all three identities $(aa)b=a(ab)$, $(ab)a=a(ba)$, $(ba)a=b(aa)$ for all
$a,b\in A$ \cite[p.~148]{Baez2002Octonions}. Any two of the three
identities imply the third, and the first two are usually taken as the
definition \cite[p.~148]{Baez2002Octonions}.

For instance, the four real division algebras $\mathbb{R},\mathbb{C},
\mathbb{H},\mathbb{O}$ are alternative; conversely, they are the only
alternative division algebras over $\mathbb{R}$ \cite[Theorem,
p.~148]{Baez2002Octonions}.

\begin{definition}[{\cite[Introduction, p.~565]{Davenport1972Malcev}}]
\label{def:nonassociative:malcev}
A \emph{Malcev algebra} $M$ over a field $F$ is a nonassociative algebra
satisfying
\begin{enumerate}
\item [(i)] $x^{2}=0$ for every $x\in M$;
\item [(ii)] $(xy)(xz) = ((xy)z)x+((yz)x)x+((zx)x)y$ for every $x,y,z\in
M$.
\end{enumerate}
\end{definition}

Malcev introduced this product as the commutator $xy-yx$ on an
alternative algebra in 1955; Sagle named the class
\cite[p.~565]{Davenport1972Malcev}. Every Lie algebra satisfies
conditions (i) and (ii), because condition (ii) then becomes the Jacobi
identity. On the trace-zero elements of the Cayley division algebra,
the commutator algebra is semisimple and seven-dimensional
\cite[p.~148]{Baez2002Octonions}, \cite[p.~52]{McCrimmon2004Taste}. % claim-ok: standard Malcev example, citations at end of sentence

Keeping the product rule of Chapter~\ref{ch:basic-theory} on the
nonassociative side, we ask the question posed and answered by the
Theorem~\ref{thm:nonassociative:davenport}.

\begin{theorem}[{\cite[Introduction and Corollary~4,
p.~565]{Davenport1972Malcev}}]
\label{thm:nonassociative:davenport}
Let $A$ be a finite-dimensional Malcev algebra over a field $F$ of
characteristic $\neq 2,3$, and suppose that the Killing form on $A$ and
the Lie multiplication algebra $L(A)$ are nondegenerate. Then every
derivation of $A$ is inner. In particular, over a field of
characteristic $0$ every derivation of a semisimple Malcev algebra is
inner.
\end{theorem}

This is the main result of \cite{Davenport1972Malcev}, proved by
reducing to the Lie algebra $L(A)$ and using the Killing-form hypothesis
\cite[p.~565]{Davenport1972Malcev}.

For the Cayley algebra over an algebraically closed field of
characteristic $0$, the derivation algebra is the exceptional simple Lie
algebra of type $G_{2}$ \cite[p.~52]{McCrimmon2004Taste}. At the group
level one has $\mathrm{Aut}(\mathbb{O})\cong G_2$
\cite[p.~159]{Baez2002Octonions}.

We can also ask what happens when we weaken the Leibniz rule itself.
Filippov studied the following class of maps for Lie algebras in earlier
work and here for alternative and Malcev algebras
\cite[Introduction, p.~354]{Filippov2000DeltaAlternative}.

\begin{definition}[{\cite[Introduction, equation~(1),
p.~354]{Filippov2000DeltaAlternative}}]
\label{def:nonassociative:delta-derivation}
Let $\Phi$ be a commutative associative ring with unity containing
$\tfrac12$, let $A$ be a $\Phi$-algebra, and let $\delta\in\Phi$ be
fixed. A \emph{$\delta$-derivation} of $A$ is a $\Phi$-linear map
$d\colon A\to A$ with
\[
d(xy)\;=\;\delta\bigl(d(x)\,y+x\,d(y)\bigr)
\qquad\text{for every }x,y\in A.
\]
Filippov calls the cases $\delta\in\{0,\tfrac12,1\}$ trivial: a
$1$-derivation is a usual derivation, a $0$-derivation is an
endomorphism vanishing on products, and every element of the centroid
of $A$ is a $\tfrac12$-derivation
\cite[Introduction, p.~354]{Filippov2000DeltaAlternative}.
\end{definition}

Filippov described $\delta$-derivations of prime alternative
$\Phi$-algebras with $\tfrac12\in\Phi$ (Theorem~\ref{thm:nonassociative:filippov-alt})
\cite[Theorem~2, p.~356]{Filippov2000DeltaAlternative}.\begin{theorem}[{\cite[Theorem~2, p.~356]{Filippov2000DeltaAlternative}}]
\label{thm:nonassociative:filippov-alt}
A prime alternative $\Phi$-algebra with $\tfrac{1}{2}\in\Phi$ has no
non-trivial $\delta$-derivations: every $\delta$-derivation is either a
derivation, a suitable multiple of the identity, or an antiderivation
in the special cases $\delta\in\{0,\tfrac12,1,-1\}$.
\end{theorem}

The content of Theorem~\ref{thm:nonassociative:filippov-malcev} is that no other maps appear. % claim-ok: own one-line summary of Theorem above, no external source

\begin{theorem}[{\cite[Theorem~3, p.~358]{Filippov2000DeltaAlternative}}]
\label{thm:nonassociative:filippov-malcev}
A prime non-Lie Malcev $\Phi$-algebra with $\tfrac{1}{2}\in\Phi$ has no
non-trivial $\delta$-derivations.
\end{theorem}

The Cayley algebra has a companion in dimension eight whose derivations are
known in every characteristic, the Okubo algebra. Let $F$ be a field and let
$(\mathcal{O},*)$ be the split Okubo algebra over $F$, a symmetric composition
algebra of dimension eight \cite[p.~4]{Elduque2015}. The derivations of
$(\mathcal{O},*)$ lie in the Lie algebra that preserves its norm (Theorem~\ref{thm:nonassociative:okubo-derivations})
\cite[p.~5]{Elduque2015}.\begin{theorem}[{\cite[Theorem~6, pp.~5--6]{Elduque2015}}]\label{thm:nonassociative:okubo-derivations}
Let $F$ be a field of characteristic not three. Then
$\operatorname{Der}(\mathcal{O},*)$ is isomorphic to $\mathfrak{sl}_3(F)$.
More precisely,
$\operatorname{Der}(\mathcal{O},*)=\{\operatorname{ad}^{*}_x\colon y\mapsto
x*y-y*x : x\in\mathcal{O}\}$.
\end{theorem}

So in characteristic not three every derivation is inner, and the underlying
vector space of $\mathcal{O}$ is $\mathfrak{sl}_3(F)$
\cite[p.~5]{Elduque2015}. In characteristic three the answer changes. The inner
derivations span a proper ideal of the derivation algebra, and the full
derivation algebra is ten-dimensional and semisimple
\cite[Theorem~9, p.~6]{Elduque2015}. Its derived algebra is simple and coincides
with the span of the inner derivations \cite[Theorem~9, p.~6]{Elduque2015}. The
derivation algebra is then isomorphic to the derivation algebra of a simple
nodal noncommutative Jordan algebra \cite[Theorem~9, p.~6]{Elduque2015}.

Section~\ref{sec:nonassociative:delta-quasi-anti} below records the
analogous $\delta$-derivation results for simple Jordan algebras and
superalgebras.

For an alternative $*$-algebra satisfying stated extra hypotheses, Pierin, Ferreira, Borges and Ferreira prove that a map is a nonlinear mixed $*$-Jordan type derivation exactly when it is an additive $*$-derivation, and draw consequences for the complex octonion algebra, associative $*$-algebras and $W^*$-factor algebras~\cite{WlAboutTheAdditivityOfANonlinear2025}.

Stitzinger proves for a finite-dimensional standard or alternative algebra over a field of characteristic zero that its derivation algebra acts completely reducibly on the algebra exactly when the radical of the algebra is contained in its centre and squares to zero~\cite{WlStandardAndAlternativeAlgebras1974}.

For a unital alternative $*$-algebra containing a nontrivial symmetric idempotent $e$ with $xA\cdot e=0$ implying $x=0$ and $xA\cdot(1_A-e)=0$ implying $x=0$, Andrade, Moraes and the two Ferreiras show that a map is a nonlinear $*$-Jordan-type derivation exactly when it is an additive $*$-derivation, and apply this to alternative $W^*$-algebras~\cite{WlNonlinearStarJordanTypeDerivat2022}.

Loos, Petersson and Racine define Lie multiplication derivations of an arbitrary nonassociative algebra over a commutative ring, describe them completely for alternative algebras following McCrimmon, and propose a new definition of inner derivations under which Schafer's standard derivations and McCrimmon's associator derivations occupy special places, showing that octonion algebras over any commutative ring have only associator derivations~\cite{WlInnerDerivationsOfAlternativeA2008}.
